\section{Results}
\label{sec:results}

\subsection{Ambiguity-Degree Probe: A Single Forward Pass Encodes Some Signal}
\label{sec:results:k-probe}

On ClarifyCodeBench restricted to $K \in \{1,2,3\}$ (excluding any
constructed $K{=}0$ class, since every tested construction of a
``no-ambiguity'' class introduced a detectable confound -- see
Section~\ref{sec:discussion}), the ambiguity-degree probe achieves test
Spearman $\rho = 0.312$ (5-fold cross-validated $\rho = 0.21 \pm 0.05$) at
its best layer (layer 32 of 37), against a random baseline of $\rho =
0.007$ and a prompt-length-only baseline of $\rho = 0.043$. The
layer-by-layer curve rises smoothly from $\rho \approx 0.08$ at layer 1 to
its layer-32 peak, consistent with a property that is progressively
constructed through the network's depth rather than a surface feature
available immediately from the embedding layer. This result is modest but
reproducible: an independent re-run with a differently-cleaned prompt
template reproduced $\rho = 0.31230908$ to eight decimal places.

Extending this probe to a pooled, four-source binary (clear-vs-ambiguous)
task spanning ClarifyCodeBench and all three Orchid sources, a
PCA-reduced logistic-regression probe reaches pooled in-distribution AUROC
of $0.75$--$0.77$; the same probe's leave-one-source-out AUROC falls to
$0.52$--$0.64$, indicating the signal generalizes only partially across
benchmark construction methodologies. A supervised-contrastive variant
trained on raw (non-PCA) activations improves leave-one-source-out AUROC to
$0.60 \pm 0.01$ (5 seeds) at some cost to pooled accuracy -- a genuine
generalization/accuracy tradeoff we revisit for the ranking probe below.

\subsection{Ranking Probe vs.\ Zero-Shot: A Cost-Quality Tradeoff, Not an Accuracy Win}
\label{sec:results:ranking}

Table~\ref{tab:ranking} reports rank-1 hit rate (did the pipeline's top
choice match the benchmark's labeled ambiguity) for the trained ranking
probe (PCA + logistic regression, the best-performing configuration) and
for zero-shot self-rank, both evaluated on the identical 1,229-group,
four-source pool under leave-one-source-out splitting.

\begin{table}[htbp]
\caption{Ranking probe vs.\ zero-shot self-rank, rank-1 hit rate, leave-one-source-out evaluation (1,229 groups).}
\label{tab:ranking}
\centering
\begin{tabular}{lcc}
\toprule
Evaluation & Probe (PCA+LR) & Zero-shot Qwen \\
\midrule
LODO mean (4 folds) & 0.637 & \textbf{0.710} \\
\quad held-out ClarifyCodeBench ($n{=}187$) & \textbf{0.701} & 0.663 \\
\quad held-out Orchid-HEval ($n{=}187$) & 0.738 & 0.743 \\
\quad held-out Orchid-BCB ($n{=}110$) & 0.582 & \textbf{0.691} \\
\quad held-out Orchid-BCB-Expand ($n{=}745$) & 0.525 & \textbf{0.742} \\
Pooled, proportional split ($n{=}246$) & 0.581 & \textbf{0.679} \\
Pooled, equal-per-source ($n{=}94$) & 0.638 & \textbf{0.713} \\
\bottomrule
\end{tabular}
\end{table}

The trained probe does not uniformly beat zero-shot self-rank: zero-shot
wins on 6 of the 7 rows in Table~\ref{tab:ranking}, most substantially on
Orchid-BCB-Expand ($0.742$ vs.\ $0.525$, a 22-point gap). The probe's one
clear win is on ClarifyCodeBench, its closest-to-training-distribution
source. We note this result explicitly reverses an earlier, smaller-scale
finding in our own development process, where the probe appeared to beat
zero-shot by 13 points on a looser, 121-group evaluation split; that
earlier result did not survive evaluation on the larger, more stringently
filtered pool used here (Section~\ref{sec:discussion}).

We therefore frame the ranking probe's contribution as a
\textbf{cost-quality tradeoff} rather than an accuracy win: it reaches
ranking quality within a few points of zero-shot self-rank on two of four
sources and its own training source, and trails meaningfully on the two
BigCodeBench-derived sources, whose prompts are long and free-form relative
to the size of the injected ambiguous edit (Section~\ref{sec:discussion}).

\subsection{Inference Cost}
\label{sec:results:cost}

% TODO(paper): replace with the fresh, apples-to-apples benchmark (batched
% N-candidate probe extraction vs. zero-shot generation, same hardware/groups)
% once GPU access is confirmed -- see plan Step 2. The number below is a
% preliminary measurement that compares zero-shot generation to a probe
% forward pass over the BARE PROMPT ONLY, not the full batched N-candidate
% extraction the deployed ranking pipeline actually performs, and should be
% treated as provisional until replaced.
On native GPU hardware (no CPU weight offload), a single probe forward pass
measures $0.079$s mean / $0.065$s median per problem, against $0.38$s mean
/ $0.37$s median for zero-shot generation with the full candidate list
inlined ($n{=}20$) -- a $4.8\times$ speedup. \textbf{[Placeholder pending a
fresh benchmark of the actual batched $N$-candidate extraction the deployed
ranking pipeline performs, per Section~\ref{sec:approach:rank-probe}; the
measurement above compares zero-shot generation to a single bare-prompt
forward pass, which is not quite the right comparison and should be treated
as provisional.]} We note that this measured advantage vanished entirely on
a memory-constrained GPU requiring CPU weight offload
(Section~\ref{sec:experimental_setup:compute}), underscoring that the
compute-efficiency argument for this class of method depends on deployment
hardware, not only on the underlying FLOP count.

\subsection{Does Ranking Quality Predict Downstream Correctness?}
\label{sec:results:correctness}

We evaluate the live two-round elicit-rank-clarify-code loop
(Section~\ref{sec:approach:loop}) against real, executable test suites on
244 problems spanning all four sources, using DeepSeek V3.2 as the coding
model. Table~\ref{tab:hit-vs-pass} reports the relationship between
round-1 ranking success (\texttt{hit1}: did the top-ranked candidate match
the labeled ambiguity) and round-1 fractional test-case pass rate,
restricted to the three Orchid sources (ClarifyCodeBench lacks Orchid's
one-injected-ambiguity, targeted-test design and is excluded from this
specific comparison).

\begin{table}[htbp]
\caption{Round-1 fractional pass rate by ranking outcome, Orchid sources, DeepSeek V3.2 coding model ($n{=}213$).}
\label{tab:hit-vs-pass}
\centering
\begin{tabular}{lccc}
\toprule
Source & hit1=1 & hit1=0 & Gap \\
\midrule
Orchid-BCB ($n{=}21$) & 88.6\% & 67.2\% & $+21.5$pp \\
Orchid-HEval ($n{=}47$) & 100.0\% & 97.9\% & $+2.1$pp \\
Orchid-BCB-Expand ($n{=}145$) & 77.5\% & 65.7\% & $+11.8$pp \\
\textbf{All Orchid combined} & \textbf{83.3\%} & \textbf{73.3\%} & \textbf{$+9.9$pp} \\
\bottomrule
\end{tabular}
\end{table}

The all-Orchid-combined gap of $+9.9$ percentage points is statistically
significant by Welch's $t$-test ($p = 0.035$) and by a 10,000-resample
bootstrap 95\% confidence interval on the mean difference, $[1.2, 19.1]$pp,
which excludes zero. A Mann-Whitney U test on the same data is more
conservative ($p = 0.10$); we report both rather than the more favorable
test alone. This is a materially stronger result than the same
relationship measured with the all-or-nothing pass metric on the same
data ($+7.6$pp, Fisher's exact $p = 0.31$, not significant) -- the
fractional metric recovers a real effect that the coarser binary metric
was not sensitive enough to detect at this sample size. We regard this
discrepancy as a methodological finding in its own right: whether an
agentic clarification pipeline appears to ``work'' can depend on which
correctness metric an evaluation uses, independent of anything about the
pipeline itself.

The effect is round-1-specific. Recomputed on round-1+2 (after a second
clarification round), the all-Orchid-combined gap falls to $+7.4$pp
($p = 0.11$, bootstrap CI includes zero); per-source, Orchid-BCB alone
strengthens to $+25.2$pp ($p = 0.05$, bootstrap CI $[5.7, 47.9]$pp,
significant) while Orchid-BCB-Expand weakens to $+8.0$pp (not
significant). We do not have a single clean explanation for this
round-to-round variability at the current sample size and report it as an
open question rather than force a narrative onto it.

We separately identify why the effect is markedly weaker for
Orchid-BCB-Expand than for the other sources: an independent LLM-judged
analysis of 1,952 candidate clarifying questions for this source finds
that, while the officially labeled ambiguity is functionally relevant to
correctness at essentially the same rate as every other source
($91.3\%$, vs.\ $91.8$--$92.4\%$ elsewhere), it is rarely the \emph{only}
relevant detail: on average $3.09$ other candidate clarifications are
independently functionally relevant to the same problem (vs.\ $1.48$ for
ClarifyCodeBench and $2.43$ for Orchid-HEval), and the label is the sole
relevant detail in only $7.0\%$ of rows (vs.\ $26.7\%$ and $14.1\%$
respectively). Resolving the one labeled ambiguity correctly therefore
caps, but does not determine, whether the generated code passes, because
most Orchid-BCB-Expand problems are underspecified in several independent
ways at once and the current pipeline resolves at most two (one per
round).

\subsection{Asking Any Question Beats Asking None}
\label{sec:results:any-question}

Independent of \emph{which} clarifying question is asked, we find that
asking any question at all is a robust, statistically significant
improvement over asking none. Comparing the same 244 problems with vs.\
without any clarification (a matched, paired comparison), round-1
clarification improves the mean fractional pass rate by $+4.0$pp (Wilcoxon
signed-rank $p = 0.0021$, paired bootstrap 95\% CI $[1.6, 6.6]$pp), and
round-1+2 by $+6.3$pp (Wilcoxon $p < 0.0001$, CI $[3.7, 9.1]$pp). Both
results are significant under every test we applied, in contrast to the
more marginal, metric-dependent ranking-quality effect in
Section~\ref{sec:results:correctness}. This suggests that, for this task
family, the decision to clarify at all carries more of the achievable
correctness gain than the precision of the ranking step that decides
\emph{which} clarification to make -- a finding with direct design
implications for agentic coding pipelines operating under a compute
budget that cannot afford a highly precise ranking step.

\subsection{A Cheap, Orthogonal Fix: Prompt Augmentation}
\label{sec:results:augmentation}

Separately from the ranking pipeline, we identify and fix a source of
trivial, crash-class failures unrelated to clarification quality: the
Orchid-BCB and Orchid-BCB-Expand ambiguity injections rewrite BigCodeBench's
free-text instruction prompt, which never carried parameter type
information to begin with, causing the coding model to guess incorrect
types for under-specified parameters (e.g., assuming a millisecond-epoch
integer timestamp was a \texttt{datetime} object). Reinstating bare
parameter types, requirements, and a worked example -- without
reinstating the full parameter-description text that would re-leak the
injected ambiguity -- and holding the ranking/clarification pipeline
completely fixed, on the 106 problems where the exact round-1/round-2
clarified prompt text could be reconstructed without a new GPU run:
round-1 mean fractional pass rate rises from $0.796$ to $0.863$, with the
zero-pass (complete-crash) rate falling from $2.8\%$ to $0.0\%$; round-2
rises from $0.804$ to $0.854$, zero-pass from $3.8\%$ to $0.9\%$. This
improvement is concentrated almost entirely in Orchid-BCB-Expand and is
orthogonal to the ranking-quality question in
Section~\ref{sec:results:correctness}: it changes prompt content only, with
the identical clarifying question selected in both conditions.
